\chapter{Certificación de productos y diseño aprobado}\label{ch:15}
\objetivos{Distinguir diseño de tipo, producto individual y cambio de diseño; relacionar pruebas, materiales e instrucciones de aeronavegabilidad continuada.}
\section{Qué acredita un certificado de tipo}
La certificación de tipo trata la conformidad de un diseño con una base aplicable y otras condiciones del proceso. No significa que cualquier ejemplar fabricado con un nombre comercial sea aeronavegable en todo momento. La conformidad de producción, el estado individual y el mantenimiento posterior tienen sus propias demostraciones y registros.

El diseño de tipo definido en 21.A.31 incluye dibujos y especificaciones, información de materiales y procesos necesaria para asegurar conformidad y otros elementos establecidos en el apartado. Esa mención conecta directamente con la primera parte del libro: composición, tratamientos y métodos de fabricación no son decisiones que puedan cambiarse libremente durante el mantenimiento.
\section{Base de certificación y demostración}
La base identifica los requisitos y especificaciones aplicables, junto con los elementos específicos que procedan. La demostración puede incluir análisis, ensayos, inspecciones y documentación. El plan de certificación organiza qué se demuestra, con qué medios y a qué requisito responde cada evidencia.

Un ensayo satisfactorio en una probeta no acredita automáticamente el componente completo. Debe justificarse representatividad de material, proceso, geometría, ambiente y carga. Tampoco una simulación que converge demuestra por sí sola seguridad: necesita modelos, hipótesis, verificación numérica y validación física proporcionadas a su función.
\ilustra{certificacion}{Cadena conceptual requisito--medio de cumplimiento--evidencia--configuración. El esquema no representa por sí solo un procedimiento de certificación completo.}
\section{Certificados suplementarios, modificaciones y reparaciones}
Un certificado de tipo suplementario (STC) aprueba un cambio en el diseño en las condiciones de la subparte correspondiente de Parte 21. Una modificación altera un diseño; una reparación restablece una condición tras daño o deterioro mediante un diseño de reparación. Las clasificaciones y rutas de aprobación se analizan en Parte 21, no se deciden solo por el tamaño visual de la intervención.

«Instalar una pieza que encaja» no demuestra conformidad con diseño aprobado. Compatibilidad funcional, configuración, instrucciones, limitaciones y documentación de instalación pueden cambiar. Un boletín de servicio, una modificación aprobada y una orden de trabajo no son documentos intercambiables.
\section{Motores: CS-E y piezas críticas}
CS-E contiene especificaciones y medios de cumplimiento para motores. En la enmienda 9 consultada, CS-E 25 trata instrucciones de aeronavegabilidad continuada; CS-E 510, análisis de seguridad; CS-E 515, piezas críticas del motor; y CS-E 740, ensayos de resistencia. La referencia debe vincularse a la base aplicable de un motor concreto, no usar la enmienda 9 como una aprobación retroactiva universal.

La identificación de piezas críticas relaciona consecuencias de fallo, propiedades del material, proceso y vida. Discos y otros elementos rotativos pueden necesitar control de vida y fabricación particularmente riguroso. Un modelo de Paris de un ejercicio no reemplaza la justificación de una pieza crítica ni sus limitaciones aprobadas.

Los ensayos de resistencia combinan secuencias y condiciones prescritas para evaluar comportamiento. No se resumen en «hacer funcionar el motor algunas horas». La temperatura, régimen, transitorios y desmontaje/inspección posteriores pueden formar parte de la evidencia según el requisito y plan aplicables.
\section{Aeronavegabilidad continuada desde el diseño}
El punto 21.A.7 establece obligaciones relativas a instrucciones para mantener el estándar de aeronavegabilidad durante la vida operacional. Las instrucciones y sus revisiones enlazan diseño y mantenimiento. No toda página de un manual tiene idéntica condición jurídica; las limitaciones de aeronavegabilidad y las instrucciones aplicables deben identificarse con precisión.

Los datos de vida, inspección y reparación tienen relación con los modos de fallo: fatiga, fluencia, degradación ambiental y daño accidental. La documentación debe permitir aplicar correctamente esos datos a la configuración individual. Una modificación puede introducir nuevas instrucciones y afectar el programa.
\section{Certificación medioambiental}
La certificación medioambiental trata requisitos de ruido, emisiones y otros aspectos dentro del ámbito y base aplicables. No se reduce a un certificado de prestaciones del motor ni sustituye su certificación de seguridad. La instalación en aeronave y las condiciones de evaluación pueden aportar requisitos distintos de los del motor aislado.

Conviene separar objetivo físico y demostración reglamentaria: reducir consumo puede mejorar eficiencia, pero demostrar una emisión requiere la magnitud, ciclo y procedimiento correspondientes. Una afirmación comercial de menor consumo no acredita cumplimiento de un requisito de emisiones.
\caso{M15. Cambio de material en una reparación}{Un taller propone sustituir una aleación de un soporte por otra con mayor resistencia nominal. Antes de evaluar ejecución se necesita determinar qué diseño y dato de reparación permiten el cambio, si es compatible con temperatura, fatiga, corrosión y uniones, y qué aprobación procede. «Mayor resistencia» no demuestra equivalencia. El caso termina en una necesidad de datos de diseño, no en autorización de sustitución.}
\section{Cadena de una modificación: diseño, instalación y mantenimiento}
Una modificación puede tener un diseño aprobado y, sin embargo, requerir todavía demostrar que su instalación es aplicable a un ejemplar determinado y se ha documentado correctamente. Las instrucciones de aeronavegabilidad continuada resultantes pueden modificar tareas, acceso, componentes, límites e inspecciones futuras. Los documentos de diseño y de mantenimiento responden a preguntas diferentes: \emph{qué configuración está aprobada} y \emph{cómo se conserva su condición de funcionamiento y aeronavegabilidad}.

Un certificado de tipo acredita un diseño en su ámbito; un certificado individual de aeronavegabilidad y un CRS documentan hechos distintos. Tampoco una aprobación STC convierte automáticamente en elegible cualquier pieza o revisión de manual. Si un cambio afecta un componente del motor, debe verificarse la combinación de diseño, número de serie, documentación e instrucciones aplicables.
\repaso{¿Es suficiente una anotación de instalación de un STC para cerrar un expediente de mantenimiento?}{No necesariamente: se necesitan el alcance, conformidad y registros prescritos, además de las instrucciones y condiciones operativas que correspondan.}


\section{Base de certificación y demostración}
La certificación de diseño establece una base y demuestra cumplimiento mediante análisis, ensayos, inspecciones y documentación pertinentes. Una simulación aislada no cubre el expediente completo. Condiciones de ensayo, conformidad del artículo y trazabilidad del resultado forman parte de la demostración. La aceptación del diseño no asegura que cualquier ejemplar fabricado o modificado conserve conformidad.

CS-E 25 trata instrucciones para aeronavegabilidad continuada; 510, análisis de seguridad; 515, partes críticas; 740, resistencia. Estos localizadores de la enmienda 9 orientan lectura, pero no sustituyen el texto completo ni la base de un motor. Citar la última enmienda no demuestra que forme parte de la base de un tipo existente.
\section{Tipo, ejemplar y configuración}
El certificado de tipo se refiere al diseño aprobado; el certificado de aeronavegabilidad del ejemplar tiene otra función. Modificaciones y reparaciones siguen sus rutas de Parte 21 según clasificación y régimen. La configuración real se compara con diseño y modificaciones aprobadas pertinentes. Que una pieza encaje físicamente no demuestra equivalencia.

Un STC tiene alcance, condiciones e información asociada. Puede incorporar suplementos e instrucciones adicionales. Conservar solo el manual base puede omitir tareas de una configuración modificada. El control de configuración conecta certificación inicial y mantenimiento diario.
\section{Información que continúa en servicio}
Las instrucciones de aeronavegabilidad continuada constituyen ese puente. 21.A.7 es una referencia actual relevante; conservar antiguas numeraciones sin comprobar puede inducir error. Disponibilidad y actualización de información se estudian con sus condiciones. No toda información del fabricante tiene la misma función u obligatoriedad.

Los problemas en servicio pueden llevar a acciones correctoras, modificaciones y AD. La experiencia devuelve información al diseño; análisis previo de seguridad y evidencia real se complementan.
\repaso{¿Por qué una pieza de mayor resistencia puede necesitar aprobación de reparación?}{También cambian rigidez, fatiga, geometría, uniones, procesos y compatibilidad. Resistencia nominal por sí sola no demuestra cumplimiento del diseño.}

\resumen{La certificación del diseño y el estado de una aeronave individual son demostraciones distintas. Materiales y procesos forman parte del diseño. STC, modificación y reparación tienen rutas y datos propios. CS-E orienta requisitos de motor dentro de una base; las instrucciones de aeronavegabilidad continuada conectan la certificación con el servicio.}
\fuentes{Reglamento 748/2012, 21.A.7, 21.A.20, 21.A.31 y subpartes D, E y M \cite{EU748}; CS-E 25, 510, 515 y 740 \cite{CSE9}; materiales docentes de Parte 21.}
